\begin{textsample}{NEG Dim 2 – 2009\_c5 – Score -8.00 – 2009\_c5/enade\_2009\_proc\_ger\_1f5954ceea.txt}  \label{ex:f2_neg_009}
Capítulo 5 Distribuição dos Conceitos 5.1 Panorama nacional da distribuição dos conceitos A Tabela 5.1 apresenta a quantidade e o percentual de cursos de TECNOLOGIA EM PROCESSOS GERENCIAIS participantes do ENADE/2009 que se situaram em cada uma das faixas de conceitos, \textbf{distribuídos} de acordo com a Região em que se localizam os cursos das instituições de educação superior, mantenedoras dos cursos responsáveis pela oferta.
O Gráfico 5.1 apresenta a mesma informação sob a forma de diagrama de caixa para uma melhor visualização da dispersão dos conceitos obtidos.
No diagrama de caixa, a caixa propriamente \textbf{dita} fica limitada pelos quartis.
A barra dentro da caixa em \textbf{negrito} representa a mediana.
Linhas verticais unem os quartis aos valores extremos.
Se estes valores extremos se configurarem como valores \textbf{atípicos}, estes aparecem como pontos isolados e as linhas terminam no maior ( ou menor ) valor ainda não considerado \textbf{atípico}.
Eventuais valores \textbf{atípicos} aparecem como pequenos círculos com a Categoria Administrativa correspondente.
Valores considerados muito \textbf{atípicos} foram representados por estrelas.
Um diagrama de caixa para poucos dados ou com dados muito concentrados pode ter alguns pontos \textbf{coincidentes}.
Por exemplo, se mais de metade dos cursos de uma Região receberem conceito 5, o valor extremo superior, o quartil superior e a mediana coincidirão.
Se o conjunto só contém um dado, este apareceria denotado como a barra da mediana.
Observando -se os dados da Tabela 5.1, percebe -se que, dos 162 cursos participantes, 41 não foram avaliados ( 25,3 \% ).
Dentre os 121 cursos avaliados em termos nacionais, \textbf{predominaram} os cursos que receberam conceitos 2 e 3, correspondendo respectivamente a 27,2 \% e a 30,2 \% dos avaliados na área de TECNOLOGIA EM PROCESSOS GERENCIAIS.
Pouco mais de 3 \% dos cursos foram classificados com conceito 1.
Conceitos 4 e 5 foram alocados a, respectivamente, 15,7 \% e 4,1 \% dos avaliados.
Os 162 cursos participantes estavam assim \textbf{distribuídos} : 8 da Região Norte, 24 da Região Nordeste, 82 da Região Sudeste, 39 da Região Sul e 9 da Região Centro-Oeste.
Considerando -se separadamente as cinco Regiões brasileiras ( ver também Gráfico 5.1 ), observa -se que, da Região Norte, 8 cursos participaram, dos quais 37,5 \% não foram avaliados.
Dentre os 5 avaliados, 4 obtiveram o conceito 2 e 1, o conceito 3.
Como a concentração no conceito 2 foi de 80 \%, cursos com conceitos acima deste valor aparecem como valores \textbf{atípicos}.
Na Região Nordeste, 24 cursos participaram dos quais 33,3 \% ( correspondendo a 8 cursos ) não foram avaliados.
Nenhum curso desta Região obteve conceito máximo, somente 1 curso ( 6,3 \% dos avaliados ) obteve conceito mínimo, 10 cursos o conceito 2, e os demais 5 cursos ficaram \textbf{distribuídos} nos conceitos 3 e 4 ( 18,8 \% e 12,5 \%, respectivamente ).
Na Região Sudeste, 82 cursos na área de TECNOLOGIA EM PROCESSOS GERENCIAIS participaram, dos quais 23,2 \% correspondendo a 19 cursos ficaram sem conceito.
Entre os avaliados, o grupo modal obteve conceito 3, sendo 25 cursos equivalentes a 39,7 \% do total regional avaliado.
O conceito máximo foi alcançado por 4 cursos ( 6,3 \% do total regional avaliado ), 14 cursos ( 22,2 \% do total regional avaliado ) obtiveram conceito 4 e 25 cursos ( 39,7 \% do total regional avaliado ) obtiveram conceito 3.
Na Região Sul, 39 cursos na área de TECNOLOGIA EM PROCESSOS GERENCIAIS participaram, dos quais 25,6 \% correspondendo a 10 cursos ficaram sem conceito.
Entre os avaliados, o grupo modal obteve conceito 3, sendo 13 cursos equivalentes a 33,3 \% do total regional avaliado.
O conceito máximo foi alcançado por 1 curso ( 3,4 \% do total regional avaliado ), 3 cursos ( 10,3 \% do total regional avaliado ) obtiveram conceito 4 e 10 cursos ( 34,5 \% do total regional avaliado ) obtiveram conceito 2.
Na Região Centro-Oeste, dos 9 cursos participantes apenas 1 ficou sem conceito e outro obteve o conceito mínimo.
Os demais 7 cursos avaliados receberam o conceito 3.
Como a concentração no conceito 3 representou 87,5 \%, semelhantemente ao que ocorreu para a Região Norte, cursos com conceitos abaixo deste valor aparecem como valores \textbf{atípicos}. 5.2 Conceitos por Categoria Administrativa e por Região A Tabela 5.2 apresenta a distribuição dos cursos participantes do ENADE/2009 de TECNOLOGIA EM PROCESSOS GERENCIAIS, por Categoria Administrativa, de acordo com os conceitos por eles alcançados, segundo Região.
O Gráfico 5.2 apresenta a mesma informação sob a forma de diagrama de caixa para uma melhor visualização da dispersão dos conceitos obtidos.
Somente eventuais valores \textbf{atípicos} aparecem com o \textbf{identificador} da Região na qual o curso foi \textbf{habilitado}.
De acordo com as informações apresentadas, a maioria dos 162 cursos participantes, 93,2 \%, é de cursos em instituições privadas, 1,9 \% dos cursos eram vinculados a instituições federais, 3,1 \% a instituições estaduais e os restantes 1,9 \% a instituições municipais.
Não foram avaliados, 41 cursos assim \textbf{distribuídos} : 1 em estadual e 40 em privadas.
As instituições federais foram as que apresentaram, em termos nacionais, maior proporção de cursos com conceito 5, equivalentes a 1/3 do total de cursos avaliados dessa categoria.
Os outros dois cursos avaliados receberam conceito 3 e 4.
Nenhum curso em instituição federal recebeu o conceito mínimo ou o conceito 2.
Entre as instituições estaduais, 2 receberam conceito 4 e os outros dois, receberam conceito 1 e 2, sendo que 1 curso não foi avaliado.
Os 3 cursos participantes em instituições municipais receberam conceito 2.
Os 151 cursos nas instituições privadas tiveram como conceito modal o valor 3 com uma incidência de 43,2 \% dos cursos avaliados ( correspondendo a 48 cursos ).
Somente 3 cursos ( 2,7 \% dos avaliados ) receberam o conceito mínimo, 4 cursos alcançaram o conceito máximo, e 16 ( 14,4 \% dos avaliados ), o conceito 4.
Entre os cursos em instituições privadas, 40 cursos ficaram sem conceito. outros dois, receberam conceito 1 e 2, sendo que 1 curso não foi avaliado.
Os 3 cursos participantes em instituições municipais receberam conceito 2.
Os 151 cursos nas instituições privadas tiveram como conceito modal o valor 3 com uma incidência de 43,2 \% dos cursos avaliados ( correspondendo a 48 cursos ).
Somente 3 cursos ( 2,7 \% dos avaliados ) receberam o conceito mínimo, 4 cursos alcançaram o conceito máximo, e 16 ( 14,4 \% dos avaliados ), o conceito 4.
Entre os cursos em instituições privadas, 40 cursos ficaram sem conceito.
Considerando -se a análise por Região observa -se que em todas as Regiões brasileiras, Instituições privadas constituem a maioria.
Na Região Norte, os 8 cursos participantes eram de instituições privadas.
Dos 8 cursos participantes em instituições privadas, 3 não foram avaliados.
Dentre os 5 avaliados, 4 obtiveram o conceito 2 e 1, o conceito 3.
Nenhum curso alcançou o conceito máximo, e nenhum curso foi alocado no conceito mínimo.
Na Região Nordeste, dos 24 cursos participantes, 22 eram de instituições privadas e 2 de instituições estaduais.
Nesta Região 8 cursos ficaram sem conceito, todos em instituições privadas.
Os 2 cursos em instituições estaduais receberam conceitos 1 e 2.
Dos 22 cursos de instituições privadas que foram avaliados, o conceito modal foi 2, 64,3 \%.
Nenhum curso em instituição privada desta Região obteve conceito máximo, nenhum obteve conceito mínimo e os demais 5 ficaram \textbf{distribuídos} nos conceitos 3 e 4 ( 21,4 \% e 14,3 \% dos avaliados, em cada conceito ).
Na Região Sudeste, dos 82 cursos participantes a maioria ( 91,5 \% ) eram de instituições privadas, 3 de instituições federais, 3 de estaduais e somente 1 de municipal.
Os 3 cursos de instituições federais receberam os conceitos 3, 4 e 5.
Dos 3 cursos em instituições estaduais, 1 não foi avaliado e os outros 2 receberam conceito 4.
O único curso em instituição municipal recebeu conceito 2.
Dos 75 cursos na área de TECNOLOGIA EM PROCESSOS GERENCIAIS em instituições privadas, 24,0 \% correspondendo a 18 cursos, ficaram sem conceito.
Entre os avaliados, o grupo modal obteve conceito 3, 24 cursos equivalentes a 42,1 \% do total regional avaliado neste tipo de instituição.
O conceito máximo foi alcançado por 3 cursos ( 5,3 \% do total regional avaliado neste tipo de instituição ), 11 cursos ( 19,3 \% do total regional avaliado neste tipo de instituição ) obtiveram conceito 4 e 19 cursos ( equivalentes a 33,3 \% do total regional avaliado neste tipo de instituição ) obtiveram conceito 2.
Na Região Sul, dos 39 cursos participantes a maioria ( 94,9 \% ) eram de instituições privadas e 2 de instituições municipais.
Os 2 cursos em instituições municipais receberam conceito 2.
Entre os 37 cursos em instituições privadas, 10 ficaram sem avaliação.
Dos 27 cursos avaliados na área neste tipo de instituição, o conceito modal foi 3 ( 13 cursos, correspondendo a 48,1 \% do total avaliado neste tipo de instituição ).
O conceito máximo foi alcançado por 1 curso e 2 cursos ( 7,4 \% do total regional avaliado neste tipo de instituição ) obtiveram conceito mínimo.
Na Região Centro-Oeste, os 9 cursos participantes eram todos de instituições privadas, sendo que 1 curso ficou sem conceito.
Entre os 8 cursos avaliados 7 obtiveram o conceito 3 e 1 o conceito mínimo. 5.3 Conceitos por Organização Acadêmica e por Região Na Tabela 5.3 encontra -se a distribuição dos conceitos atribuídos aos cursos participantes do ENADE/2009 de TECNOLOGIA EM PROCESSOS GERENCIAIS, por Organização Acadêmica, segundo as Regiões.
O Gráfico 5.3 apresenta a mesma informação sob a forma de diagrama de caixa para uma melhor visualização da dispersão dos conceitos obtidos.
Somente eventuais valores \textbf{atípicos} apareceram com o \textbf{identificador} da Categoria Administrativa.
Os 161 cursos desta área estão principalmente alocados em Faculdades ( 47,5 \% correspondendo a 77 cursos ), depois em Universidades ( 34,6 \% correspondendo a 56 cursos ) e Centros Universitários ( 16,0 \% correspondendo a 26 cursos ).
Este ordenamento das 3 categorias principais se repete em todas as Regiões, com exceção da Região Centro-Oeste, onde cursos de Universidades não participaram.
Participaram ainda 3 cursos em IFECT.
Dos 3 cursos em IFECT, todos foram avaliados e receberam conceitos 3, 4 e 5.
Entre os 41 cursos não avaliados, 13 estavam alocados em Universidades, 4 em Centros Universitários e 24 em Faculdades.
Universitários e 24 em Faculdades.
Nas Universidades o conceito modal foi 2 ( 46,5 \% dos avaliados em Universidades ), nas Faculdades o conceito modal ficou dividido nos valores 2 e 3 ( 37,7 \% cada ), nos Centros Universitários, o conceito modal foi 3 ( 59,1 \% ).
Poucos cursos nos 3 tipos de Organização Acadêmica mais frequentes, Universidades, Centros Universitários e Faculdades, receberam conceito mínimo, respectivamente 1, 1 e 2 cursos, respectivamente.
Somente 1 curso em Faculdade, 1 em IFECT e 3 em Universidades atingiram o conceito máximo.
Considerando -se separadamente as cinco Regiões, observa -se que, a Região Norte não apresentou nenhum curso em CEFET, nenhum em IFECT e nenhum na categoria residual.
Dos 8 cursos participantes, 1 estava alocado em Universidade, 1 em Centros Universitários e 6 em Faculdades.
O único curso avaliado em Universidade recebeu conceito 2.
O único curso participante em Centro Universitário não foi avaliado.
Dos 6 cursos em Faculdades, dois não foram avaliados e 1 recebeu o conceito 3 e 3 o conceito 2.
Nenhum curso nesta Região obteve nem o conceito máximo, nem o conceito mínimo e nem o conceito 4.
A Região Nordeste não apresentou nenhum curso na área de TECNOLOGIA EM PROCESSOS GERENCIAIS em CEFET ou IFECT ou na categoria residual.
Dos 24 cursos, 8 cursos não foram avaliados, todos em Faculdades.
Os demais 16 cursos estavam assim \textbf{distribuídos} : 4 em Universidades, 1 em Centro Universitário e 11 em Faculdades.
Dos 4 cursos em Universidades na Região, 1 recebeu o conceito mínimo e 3, o conceito 2.
O único curso em Centro Universitário recebeu o conceito 4.
Entre os 11 cursos em Faculdades na Região,. o conceito modal foi 2, com 7 cursos, seguido do conceito 3, com 3 cursos e o conceito 4, com somente 1 curso.
Nenhum curso nesta Região obteve conceito máximo.
A Região Sudeste apresentou 82 cursos na área de TECNOLOGIA EM PROCESSOS GERENCIAIS assim \textbf{distribuídos} : 36 cursos nas Universidades, 16 nos Centros Universitários, 27 nas Faculdades e 3 cursos em IFECT.
Dos 36 cursos alocados em Universidades, 9 não foram avaliados.
Os cursos em Universidades, receberam conceitos 2, 3, 4 e 5 ( respectivamente 11, 10, 4 e 2 cursos ).
Nos Centros Universitários 2 cursos não foram avaliados e os demais receberam os conceitos 2, 3 e 4 ( respectivamente 3, 8 e 3 cursos ).
Dos 27 cursos em Faculdades na Região, 8 não foram avaliados.
Os 19 restantes dividiram -se igualmente entre os conceitos 2, 3 e 4 com 6 cursos cada ( 31,6 \% dos avaliados neste tipo de instituição na Região ).
Apenas 1 curso recebeu o conceito máximo em Faculdades.
A Região Sul apresentou curso na área de TECNOLOGIA EM PROCESSOS GERENCIAIS somente nas 3 categorias mais frequentes : 15 em Universidades ( 38,5 \% ), 4 em Centros Universitários ( 10,3 \% ) e 20 em Faculdades ( 51,3 \% ). avaliados, 4 estavam em Universidades, 1 em Centro Universitário e 5 em Faculdades.
Os 11 cursos avaliados em Universidades da Região receberam conceitos 2, 3 e 5 ( respectivamente 5, 5 e 1 curso ).
Os 3 cursos avaliados em Centros Universitários receberam conceitos 1, 2 e 3.
Entre os cursos avaliados em Faculdades, o conceito modal foi o conceito 3 ( 7 cursos correspondendo a 46,7 \% dos cursos ).
Os demais cursos receberam conceitos 1, 2 e 4 ( respectivamente, 1, 4 e 3 cursos ).
A Região Centro-Oeste só apresentou cursos na área de TECNOLOGIA EM PROCESSOS GERENCIAIS em 2 categorias mais frequentes : Centros Universitários ( 4 cursos ) e Faculdades ( 5 cursos ).
Os 4 cursos em Centros Universitários receberam conceito 3.
Dos 5 cursos em Faculdades, 3 receberam o conceito 3, 1 o conceito mínimo e outro não recebeu conceito.

% matched lemmas: atípico, coincidente, distribuir, dito, habilitado, identificador, negrito, predominar
\end{textsample}
